% GENERATED FILE. Edit canonical lessons, then run npm run generate:books.
% canonical-source-hash: fnv1a64:4d9950e72dc1b363
% canonical-lessons: SA-C403-panthah, SA-C403-dirghah, SA-C403-durgamah, SA-C403-padatih, SA-C403-vicarati

\chapter{The Long Road}
\label{ch:sa-the-long-road}
\hlchaptermodality{403}

\begin{chapteropening}
\textbf{By the end of this chapter:} \emph{I can describe a long, hard road and who walks it.}

Complete the last lesson of chapter 403: I can describe a long, hard road and who walks it.
\end{chapteropening}

\section[\texorpdfstring{\sk{पन्थाः} (panthāḥ)}{पन्थाः (panthāḥ)}]{\sk{पन्थाः} (panthāḥ) — a route; a path}

\label{lesson:SA-C403-panthah}



\begin{quote}
Before the new one: say the Sanskrit for lies down to sleep, then the Sanskrit for water for washing the feet.
\end{quote}

\subsection*{You'll want to know: \sk{पन्थाः}}
\textbf{\sk{पन्थाः}} — \emph{panthāḥ} — \textquotedblleft{}a route; a path\textquotedblright{}.

\begin{quote}
\sk{पन्थाः} \sk{वनस्य} \sk{मध्ये} \sk{गच्छति।}
\end{quote}

\begin{quote}
\emph{panthāḥ vanasya madhye gacchati.}
\end{quote}

\begin{quote}
The route runs through the middle of the forest.
\end{quote}

A masculine noun.

\begin{cousinweb}
\textbf{\sk{पन्थाः}} (\emph{panthāḥ}) is the nominative of \emph{pathin}, \textquotedblleft{}a way, a road\textquotedblright{}, whose strong stem is \emph{panthan}; it comes from the same word as \emph{pātheya}, \textquotedblleft{}food for the journey\textquotedblright{}; MW gives \emph{pathin} as a way, path, road or route (MW).
\end{cousinweb}

\begin{grammarlens}[title={What you've built}]
One more word for getting around.
\end{grammarlens}

\subsection*{Your turn}
\begin{itemize}
\raggedright
  \item \emph{Say it:} \emph{panthāḥ}
  \item \emph{Say it:} \emph{panthāḥ}, once more
  \item \emph{Say it:} \emph{panthāḥ vanasya madhye gacchati.} — The route runs through the middle of the forest.
  \item \emph{Recall:} say the Sanskrit for lies down to sleep, then the Sanskrit for water for washing the feet, then say \emph{panthāḥ} again
\end{itemize}

\begin{tcolorbox}[breakable,colback=teal!4,colframe=teal!35!black,title={Before you move on}]
What does \sk{पन्थाः} mean? (\textquotedblleft{}A route; a path\textquotedblright{}.) Say it once more.
\end{tcolorbox}
\section[\texorpdfstring{\sk{दीर्घः} (dīrghaḥ)}{दीर्घः (dīrghaḥ)}]{\sk{दीर्घः} (dīrghaḥ) — long}

\label{lesson:SA-C403-dirghah}



\begin{quote}
Before the new one: say the Sanskrit for water for washing the feet, then the Sanskrit for a route.
\end{quote}

\subsection*{You'll want to know: \sk{दीर्घः}}
\textbf{\sk{दीर्घः}} — \emph{dīrghaḥ} — \textquotedblleft{}long\textquotedblright{}.

\begin{quote}
\sk{पर्वतं} \sk{प्रति} \sk{पन्थाः} \sk{दीर्घः} \sk{अस्ति।}
\end{quote}

\begin{quote}
\emph{parvataṃ prati panthāḥ dīrghaḥ asti.}
\end{quote}

\begin{quote}
The route to the mountain is long.
\end{quote}

\begin{cousinweb}
\textbf{\sk{दीर्घः}} (\emph{dīrghaḥ}) is the adjective MW gives for long, in space and in time; MW gives no derivation, so its origin is uncertain (MW).
\end{cousinweb}

\begin{grammarlens}[title={What you've built}]
One more word for getting around.
\end{grammarlens}

\subsection*{Your turn}
\begin{itemize}
\raggedright
  \item \emph{Say it:} \emph{dīrghaḥ}
  \item \emph{Say it:} \emph{dīrghaḥ}, once more
  \item \emph{Say it:} \emph{parvataṃ prati panthāḥ dīrghaḥ asti.} — The route to the mountain is long.
  \item \emph{Recall:} say the Sanskrit for water for washing the feet, then the Sanskrit for a route, then say \emph{dīrghaḥ} again
\end{itemize}

\begin{tcolorbox}[breakable,colback=teal!4,colframe=teal!35!black,title={Before you move on}]
What does \sk{दीर्घः} mean? (\textquotedblleft{}Long\textquotedblright{}.) Say it once more.
\end{tcolorbox}
\section[\texorpdfstring{\sk{दुर्गमः} (durgamaḥ)}{दुर्गमः (durgamaḥ)}]{\sk{दुर्गमः} (durgamaḥ) — hard to get through; impassable}

\label{lesson:SA-C403-durgamah}



\begin{quote}
Before the new one: say the Sanskrit for a route, then the Sanskrit for long.
\end{quote}

\subsection*{You'll want to know: \sk{दुर्गमः}}
\textbf{\sk{दुर्गमः}} — \emph{durgamaḥ} — \textquotedblleft{}hard to get through; impassable\textquotedblright{}.

\begin{quote}
\sk{वर्षायां} \sk{पन्थाः} \sk{दुर्गमः} \sk{भवति।}
\end{quote}

\begin{quote}
\emph{varṣāyāṃ panthāḥ durgamaḥ bhavati.}
\end{quote}

\begin{quote}
In the rains the route becomes hard to get through.
\end{quote}

\begin{cousinweb}
\textbf{\sk{दुर्गमः}} (\emph{durgamaḥ}) is built from \emph{dur}, \textquotedblleft{}hard, ill\textquotedblright{}, and the root √gam, \textquotedblleft{}to go\textquotedblright{}: hard to go through; MW gives \emph{durgama} as difficult to be traversed, impassable (MW).
\end{cousinweb}

\begin{grammarlens}[title={What you've built}]
One more word for getting around.
\end{grammarlens}

\subsection*{Your turn}
\begin{itemize}
\raggedright
  \item \emph{Say it:} \emph{durgamaḥ}
  \item \emph{Say it:} \emph{durgamaḥ}, once more
  \item \emph{Say it:} \emph{varṣāyāṃ panthāḥ durgamaḥ bhavati.} — In the rains the route becomes hard to get through.
  \item \emph{Recall:} say the Sanskrit for a route, then the Sanskrit for long, then say \emph{durgamaḥ} again
\end{itemize}

\begin{tcolorbox}[breakable,colback=teal!4,colframe=teal!35!black,title={Before you move on}]
What does \sk{दुर्गमः} mean? (\textquotedblleft{}Hard to get through; impassable\textquotedblright{}.) Say it once more.
\end{tcolorbox}
\section[\texorpdfstring{\sk{पदातिः} (padātiḥ)}{पदातिः (padātiḥ)}]{\sk{पदातिः} (padātiḥ) — a traveller on foot}

\label{lesson:SA-C403-padatih}



\begin{quote}
Before the new one: say the Sanskrit for long, then the Sanskrit for hard to get through.
\end{quote}

\subsection*{You'll want to know: \sk{पदातिः}}
\textbf{\sk{पदातिः}} — \emph{padātiḥ} — \textquotedblleft{}a traveller on foot\textquotedblright{}.

\begin{quote}
\sk{पदातिः} \sk{मार्गस्य} \sk{पार्श्वे} \sk{चलति।}
\end{quote}

\begin{quote}
\emph{padātiḥ mārgasya pārśve calati.}
\end{quote}

\begin{quote}
The traveller on foot walks at the side of the road.
\end{quote}

A masculine noun.

\begin{cousinweb}
\textbf{\sk{पदातिः}} (\emph{padātiḥ}) is built on \emph{pada}, \textquotedblleft{}a foot, a step\textquotedblright{}; MW gives \emph{padāti} as going on foot, a pedestrian, and divides it, with a query, as \emph{pada} + \emph{āti} (MW).
\end{cousinweb}

\begin{grammarlens}[title={What you've built}]
One more word for getting around.
\end{grammarlens}

\subsection*{Your turn}
\begin{itemize}
\raggedright
  \item \emph{Say it:} \emph{padātiḥ}
  \item \emph{Say it:} \emph{padātiḥ}, once more
  \item \emph{Say it:} \emph{padātiḥ mārgasya pārśve calati.} — The traveller on foot walks at the side of the road.
  \item \emph{Recall:} say the Sanskrit for long, then the Sanskrit for hard to get through, then say \emph{padātiḥ} again
\end{itemize}

\begin{tcolorbox}[breakable,colback=teal!4,colframe=teal!35!black,title={Before you move on}]
What does \sk{पदातिः} mean? (\textquotedblleft{}A traveller on foot\textquotedblright{}.) Say it once more.
\end{tcolorbox}
\section[\texorpdfstring{\sk{विचरति} (vicarati)}{विचरति (vicarati)}]{\sk{विचरति} (vicarati) — roams; rambles about}

\label{lesson:SA-C403-vicarati}



\begin{quote}
Before the new one: say the Sanskrit for hard to get through, then the Sanskrit for a traveller on foot.
\end{quote}

\subsection*{You'll want to know: \sk{विचरति}}
\textbf{\sk{विचरति}} — \emph{vicarati} — \textquotedblleft{}roams; rambles about\textquotedblright{}.

\begin{quote}
\sk{पथिकः} \sk{वने} \sk{विचरति।}
\end{quote}

\begin{quote}
\emph{pathikaḥ vane vicarati.}
\end{quote}

\begin{quote}
The wayfarer roams in the forest.
\end{quote}

Past (he …): \sk{सः} \sk{वने} \sk{विचरितवान्}, \emph{saḥ vane vicaritavān}.

\begin{cousinweb}
\textbf{\sk{विचरति}} (\emph{vicarati}) is built from \emph{vi}, \textquotedblleft{}apart, here and there\textquotedblright{}, and the root √car, \textquotedblleft{}to move\textquotedblright{}, which also gives \emph{carati}, \textquotedblleft{}moves about\textquotedblright{}; MW gives \emph{vi-car} as to rove or ramble about (MW).
\end{cousinweb}

\begin{grammarlens}[title={What you've built}]
One more word for getting around.
\end{grammarlens}

\subsection*{Your turn}
\begin{itemize}
\raggedright
  \item \emph{Say it:} \emph{vicarati}
  \item \emph{Say it:} \emph{vicarati}, once more
  \item \emph{Say it:} \emph{pathikaḥ vane vicarati.} — The wayfarer roams in the forest.
  \item \emph{Recall:} say the Sanskrit for hard to get through, then the Sanskrit for a traveller on foot, then say \emph{vicarati} again
\end{itemize}

\begin{tcolorbox}[breakable,colback=teal!4,colframe=teal!35!black,title={Before you move on}]
What does \sk{विचरति} mean? (\textquotedblleft{}Roams; rambles about\textquotedblright{}.) Say it once more.
\end{tcolorbox}
